\documentclass[11pt,a4paper]{article}

\usepackage[margin=1in]{geometry}
\usepackage{lmodern}
\usepackage[T1]{fontenc}
\usepackage{amsmath,amssymb}
\usepackage{booktabs}
\usepackage{longtable}
\usepackage{xcolor}
\usepackage{listings}
\usepackage{enumitem}
\usepackage{fancyhdr}
\usepackage[hidelinks,colorlinks=true,linkcolor=blue!55!black,urlcolor=blue!55!black,citecolor=blue!55!black]{hyperref}

% ---- literal underscores etc. in text ----
\usepackage[strings]{underscore}

% ---- breakable DOI links (must follow hyperref) ----
\usepackage{xurl}
\newcommand{\doi}[1]{\href{https://doi.org/#1}{\nolinkurl{#1}}}

% ---- colours ----
\definecolor{codebg}{RGB}{245,245,248}
\definecolor{codekw}{RGB}{0,80,160}
\definecolor{codecm}{RGB}{0,128,64}
\definecolor{codestr}{RGB}{160,60,20}
\definecolor{rule}{RGB}{40,60,120}

% ---- listings ----
\lstdefinestyle{sh}{
  basicstyle=\ttfamily\small,
  backgroundcolor=\color{codebg},
  frame=single, rulecolor=\color{rule!30},
  breaklines=true, columns=fullflexible,
  keepspaces=true, showstringspaces=false,
  keywordstyle=\color{codekw}\bfseries,
  commentstyle=\color{codecm}\itshape,
  stringstyle=\color{codestr},
  xleftmargin=8pt, xrightmargin=4pt, aboveskip=6pt, belowskip=6pt
}
\lstdefinestyle{mat}{style=sh,
  morekeywords={function,end,for,while,if,else,elseif,global,parfor,parfeval,run,addpath,clear}
}
\lstset{style=sh}
\newcommand{\code}[1]{\texttt{#1}}

% ---- headers ----
\pagestyle{fancy}\fancyhf{}
\lhead{\small HEBCPF -- MATLAB v3.202606}
\rhead{\small User Guide}
\cfoot{\small\thepage}
\renewcommand{\headrulewidth}{0.4pt}

\setlength{\parskip}{4pt}
\setlength{\parindent}{0pt}

\begin{document}

% ================= TITLE =================
\begin{titlepage}
\centering
\vspace*{2.2cm}
{\color{rule}\rule{\textwidth}{1.2pt}}\\[8pt]
{\Huge\bfseries HEBCPF}\\[10pt]
{\LARGE Holomorphic-Embedding-Based Continuation}\\[4pt]
{\LARGE Power-Flow All-Solutions Solver}\\[10pt]
{\color{rule}\rule{\textwidth}{1.2pt}}\\[18pt]
{\large\itshape Finding \emph{all} real-valued solutions of the AC power-flow equations}\\[6pt]
{\large\itshape via the ellipsoidal formulation, holomorphic embedding, Pad\'e approximation, and numerical continuation}\\[26pt]

{\large\bfseries Release: MATLAB v3.202606}\\[4pt]
{\large Pure-MATLAB build (no compiled binaries; cross-platform) --- validated, robustness-hardened}\\[30pt]

\begin{tabular}{rl}
\textbf{Method:} & Holomorphic Embedding Based Continuation (HEBC)\\[3pt]
\textbf{Originating work:} & D.\ Wu and B.\ Wang, \emph{IEEE Access}, 2019\\[6pt]
\textbf{Foundations:} & Elliptical branch tracing --- B.\ Lesieutre \& D.\ Wu,\\[1pt]
                      & Allerton 2015;\\[3pt]
                      & Algebraic set preserving mappings --- D.\ Wu,\\[1pt]
                      & Ph.D.\ dissertation, UW--Madison, 2017\\[6pt]
\textbf{Document:} & User Guide and Reference Manual\\
\end{tabular}

\vfill
{\small When you publish results obtained with this solver, please cite the HEBC
paper (see Section~\ref{sec:cite}).}\\[6pt]
{\small\today}
\end{titlepage}

\tableofcontents
\newpage

% ================= 1. OVERVIEW =================
\section{Overview}

HEBCPF computes \textbf{all} real-valued solutions of the AC power-flow (load-flow)
equations for a given network, rather than the single operating point returned by a
conventional solver. Knowing the complete solution set is essential for
transient-stability energy-function methods (locating type-1 unstable equilibria),
voltage-stability analysis, and the study of multiple high-voltage operating points.

The solver implements the \textbf{Holomorphic Embedding Based Continuation (HEBC)}
method of Wu and Wang (see Section~\ref{sec:cite}). The power-flow equations are first
recast in an \emph{elliptical form} (Lesieutre and Wu, 2015) so that every
solution-tracing curve is a
bounded loop; the loops are then followed by a hybrid numerical scheme:

\begin{itemize}[leftmargin=1.4em]
  \item \textbf{Holomorphic embedding} traverses the smooth (non-singular) portions of
        a curve in a few very large steps, using power-series / Pad\'e evaluation of the
        holomorphically-embedded voltages;
  \item a classical \textbf{predictor--corrector continuation} takes over only in the
        neighbourhood of singular points (turning points / folds), where it steps
        carefully through the fold and hands control back to the holomorphic routine.
\end{itemize}

Each time the tracing (homotopy) parameter crosses zero, a power-flow solution has been
found. Starting from one known operating point (from MATPOWER's \code{runpf}), the solver
recursively traces the curves emanating from every discovered solution until no new
solution appears.

This release adds a \textbf{consensus Pad\'e-pole estimator} to the classic method
(Section~\ref{sec:pole}), which stabilises the placement of the holomorphic/continuation
hand-off by agreeing a singularity estimate across several buses rather than trusting a
single bus.

Three execution frameworks share the identical numerical hot path:

\begin{center}
\begin{tabular}{@{}lll@{}}
\toprule
Framework & Entry script & Best for\\
\midrule
Serial            & \code{run_batch.m}          & small cases, debugging, reproducibility\\
Parallel (parfor) & \code{run_batch_par.m}      & medium cases, one worker per equation\\
Parallel (parfeval)& \code{run_batch_parfeval.m} & \textbf{large cases (recommended)}\\
\bottomrule
\end{tabular}
\end{center}

% ================= 2. REQUIREMENTS =================
\section{Requirements and dependencies}

\begin{description}[leftmargin=0em,style=nextline]
\item[MATLAB (required).] MATLAB R2022a or later, on \emph{any} platform --- this is a
      pure-MATLAB build with no compiled binaries. Uses \code{sparse}, \code{linsolve},
      \code{lu}, and standard linear algebra.
\item[MATPOWER (required).] MATPOWER must be installed and on the MATLAB path. The setup
      stage calls MATPOWER functions including \code{runpf} for the initial operating
      point, \code{makeYbus} for the bus-admittance matrix, and
      \code{idx_bus}/\code{idx_brch} for MATPOWER matrix indices. The solver reads the
      bundled MATPOWER-format \code{caseXXX.m} data files, but does \emph{not} ship
      \code{runpf}, \code{makeYbus}, \code{idx_bus}, \code{idx_brch}, or their
      dependencies --- those come from MATPOWER (\url{https://matpower.org}).
\item[Parallel Computing Toolbox (optional).] Required \emph{only} for the two parallel
      frameworks (\code{run_batch_par.m}, \code{run_batch_parfeval.m}). The serial
      framework runs without it.
\item[No compiler needed.] This is the pure-MATLAB build: it runs on any MATLAB platform
      with no MEX / codegen step. For roughly $2\times$ faster serial tracing on Windows
      x64, use the companion compiled build (\code{HEBCPF_MEX_v3.202606}); the numerical
      results are identical.
\end{description}

% ================= 3. INSTALLATION =================
\section{Installation}

\begin{enumerate}[leftmargin=1.6em]
  \item Install MATLAB and MATPOWER; confirm \code{runpf} works
        (\code{which runpf} returns a path).
  \item Unpack this folder (\code{HEBCPF_matlab_v3.202606}) anywhere on disk.
  \item Add the folder to the MATLAB path each session:
\begin{lstlisting}
>> cd '<path>\HEBCPF_matlab_v3.202606'
>> addpath(pwd)
\end{lstlisting}
  \item Add only this one HEBCPF release folder to the MATLAB path. Do not add all four
        release folders simultaneously, because they contain many functions and cases with
        the same names.
  \item Nothing else is needed --- this build is pure MATLAB, so there are no binaries to
        install or compile, on any platform.
\end{enumerate}

% ================= 4. QUICK START =================
\section{Quick start}

Before running any driver, confirm MATLAB can find MATPOWER:
\begin{lstlisting}[style=mat]
>> which runpf
>> which makeYbus
\end{lstlisting}
Both commands should return paths inside the MATPOWER installation.

From a system shell (non-interactive \code{-batch}):
\begin{lstlisting}[style=sh]
:: SERIAL
matlab -sd "<path>\HEBCPF_matlab_v3.202606" -batch "addpath(pwd); run('run_batch.m')"

:: PARALLEL (parfor)
matlab -sd "<path>\HEBCPF_matlab_v3.202606" -batch "addpath(pwd); run('run_batch_par.m')"

:: PARALLEL (parfeval, recommended for big cases)
matlab -sd "<path>\HEBCPF_matlab_v3.202606" -batch "addpath(pwd); run('run_batch_parfeval.m')"
\end{lstlisting}

From an interactive MATLAB session:
\begin{lstlisting}[style=mat]
>> addpath(pwd)
>> run_batch            % serial batch
>> run_batch_par        % parallel (parfor)
>> run_batch_parfeval   % parallel (parfeval)
>> r = run_merged_case('case14mod','parfeval');  % one case, any mode
>> r.solutions, r.max_residual, r.wall_time
\end{lstlisting}

The two parallel drivers start a local \code{parpool} automatically (default worker
count) if none is running. The one-time pool start-up is excluded from the reported
timing.

% ================= 5. USER OPTIONS =================
\section{User options}

\subsection{Where the settings live: \code{solver_params.m}}
Every tunable constant is defined in one place --- the function \code{solver_params.m},
which returns a \code{param} struct consumed by the whole solver (\code{hybrid_traceloops}
$\to$ holomorphic step, continuation, corrector). To change any behaviour, edit the
relevant line in \code{solver_params.m}; all three frameworks \emph{and} the parallel
workers read the same file, so one edit applies everywhere. Every field is documented
inline, and Section~\ref{sec:params} tables the ones you are most likely to touch. (The
only tunable kept outside this file is the Pad\'e \code{degree}, set in the run drivers and
\code{main.m}.)

\subsection{Choosing which cases to run}
Each driver has a list near the top:
\begin{lstlisting}[style=mat]
cases_to_run = { 'case9', 'case14mod', 'case39', ... };
\end{lstlisting}
Comment out or add case names. Bundled test systems (expected solution counts, collapsing
antipodal pairs, in parentheses):

\begin{center}\small
\begin{tabular}{@{}llll@{}}
\toprule
\code{case3} (6) & \code{case4gs} (6) & \code{case4BB0} (14) & \code{case4BBc} (12)\\
\code{case5loop} (10) & \code{case5Salam} (10) & \code{case6ww} (6) & \code{case7Salam} (4)\\
\code{case9} (8) & \code{case9Q} (8) & \code{case14mod} (30) & \code{case33bw} (16)\\
\code{case39} (176) & \code{case30}/\code{case_ieee30} (472) & \code{case57mod} (606) & \code{case57} (1322)\\
\bottomrule
\end{tabular}
\end{center}

\subsection{Load scaling}
Each driver contains
\begin{lstlisting}[style=mat]
factor = 1 + 0.19*(0);   % default = 1 (nominal load)
\end{lstlisting}
Changing the multiplier (the ``\code{(0)}'') scales all bus loads and generator outputs,
letting you sweep load level. Leave it at 1 to reproduce the nominal solution counts.

\subsection{Choosing a framework: serial, \code{parfor}, \code{parfeval}}\label{sec:frameworks}
All three frameworks compute the \emph{identical} solution set; they differ only in how the
per-equation curve traces are distributed across processes.
\begin{description}[leftmargin=0em,style=nextline]
\item[Serial (\code{run_batch.m})] one process traces every curve in turn. Needs no
      toolbox; fully deterministic and the easiest to debug.
\item[\code{parfor} (\code{run_batch_par.m})] MATLAB's parallel \emph{for}-loop. The
      equations of one starting solution are grouped into a \emph{batch} and dealt out to
      the worker pool, but the loop then \textbf{blocks until the whole batch has
      finished} before it proceeds. It is simple and effective, yet a batch runs only as
      fast as its \emph{slowest} equation: if one curve is long or stalls, the other
      workers sit idle at that synchronisation barrier until it completes.
\item[\code{parfeval} (\code{run_batch_parfeval.m})] \emph{asynchronous} task submission.
      Each equation is submitted individually via \code{parfeval}, and results are
      collected in the order they \emph{complete}, whatever that order is. A sliding window
      keeps the pool saturated: the instant a worker finishes a curve it is handed the next
      one, \emph{without} waiting for the others. This removes the per-batch barrier and
      keeps every worker busy, so it is the fastest framework on large systems with uneven
      curve lengths (e.g.\ \code{case30}, \code{case57mod}). The only trade-off is that the
      \emph{order} in which curves are traced is non-deterministic (it depends on which
      worker finishes first); the final solution set is identical, but which curve happens
      to discover a given solution can differ between runs.
\end{description}
In one line: \code{parfor} is \emph{synchronous, batch-at-a-time} and blocks on the slowest
curve; \code{parfeval} is \emph{asynchronous, task-at-a-time} with no idle waiting ---
prefer it for large cases. Both start a local pool automatically and require the Parallel
Computing Toolbox; the one-time pool start-up is excluded from the reported timing.

\subsection{Pad\'e pole estimation (consensus)}\label{sec:pole}
The hand-off from the holomorphic routine to the continuation routine is driven by an
estimate of the nearest singularity, taken from the real roots (poles) of the Pad\'e
denominator. This build estimates that pole by \textbf{consensus} across a fixed set of
buses (\code{param.pole_buses}, default \code{2:6}, auto-filtered to $\le$ bus count): the
true singularity sits at nearly the same magnitude on every bus, whereas spurious
(Froissart) poles are bus-specific. The estimator takes the smallest pole magnitude agreed
on by at least half of the sampled buses (clustered within \code{param.pole_cl_tol}), and
falls back to the global smallest if no consensus forms. This replaces the legacy
single-/random-bus pole pick and is one of the robustness improvements in this build
(Section~\ref{sec:tuned}). To change it, edit \code{param.pole_buses},
\code{param.pole_consensus}, or \code{param.pole_cl_tol} in \code{solver_params.m}.

\subsection{Programmatic single-case runs: \code{run_merged_case}}
While the \code{run_batch*} drivers loop over a list of cases and write CSV summaries,
\code{run_merged_case} is the \textbf{unified single-case harness}: it runs \emph{one}
case in \emph{any} of the three frameworks behind a single function call and returns a
structured, machine-readable result instead of printing to the console.
\begin{lstlisting}[style=mat]
[result, solutions] = run_merged_case(case_name, mode, close_pool_on_finish)
%   mode = 'serial' | 'parfor' | 'parfeval'
\end{lstlisting}
It performs the identical preprocessing as the batch drivers (it reuses \code{main.m} as
the setup template, substituting the requested case for the active \code{mpc=case...}
line), traces the case in the selected framework, and \textbf{self-verifies} by
recomputing the maximum residual $|z^\top M_i z - b_i|$ over every solution $z$ and every
equation $i$. The call never throws: on failure it returns a \code{FAIL} result rather
than erroring. The returned \code{result} struct contains:

\begin{center}\small
\begin{tabular}{@{}ll@{}}
\toprule
Field & Meaning\\
\midrule
\code{status} & \code{'PASS'} or \code{'FAIL'}\\
\code{case_name}, \code{mode} & echoed inputs\\
\code{buses} & number of buses\\
\code{solutions} & number of solutions found\\
\code{max_residual} & largest constraint residual over the solution set (correctness check)\\
\code{wall_time} & trace wall-clock time (excludes pool start-up)\\
\code{error_id}, \code{message} & populated only on \code{FAIL}\\
\bottomrule
\end{tabular}
\end{center}
The second output, \code{solutions}, is the full solution matrix \code{Zsave} (each column
a solution). Typical uses: A/B comparison of the three frameworks on one case (identical
setup, compare \code{wall_time}/\code{solutions}/\code{max_residual}); automated
regression or validation checks; and obtaining the solution set and residual directly as
return values. Because it reads \code{main.m}, that file must remain in the folder for
this harness to work.

% ================= 6. HOW THE SOLVER WORKS =================
\section{How the solver works}

\subsection{Elliptical formulation}
In rectangular voltage coordinates $x=[V_d^\top\ V_q^\top]^\top$, each power-flow
equation is a quadratic $x^\top C_i x = b_i$. These native surfaces are unbounded
(hyperbolic), so raw branch tracing can run off to infinity. Following
Lesieutre--Wu (2015) and Wu's thesis, a positive-definite \emph{base
ellipsoid} is added to every equation, producing an equivalent system
\[
  x^\top M_i x = 1, \qquad M_i = M_i^\top \succ 0, \quad i=1,\dots,2N_{\text{bus}},
\]
whose zero set is \emph{identical} to the original (algebraic-set preserving) but whose
every equation is now a high-dimensional ellipse centred at the origin. Relaxing one
equation with a tracing parameter $\alpha$ (the code variable \code{aal}) yields a
1-dimensional curve; because the ellipses are compact, each connected component of that
curve is a \textbf{bounded loop} (generically diffeomorphic to a circle), so a trace must
return to its start. Solutions occur wherever $\alpha=0$ along the loop.

\subsection{Holomorphic embedding and Pad\'e}
Rather than crawl the loop with small predictor--corrector steps, HEBC embeds the tracing
parameter into the complex plane and represents each bus voltage as a holomorphic
function, computed as a power series (to degree \code{degree.act}, default 15). All
series degrees share a single LU factorisation, so the coefficients are cheap. A
\textbf{Pad\'e approximant} (balanced numerator/denominator degree, for maximal
convergence radius) then extrapolates the voltages far along the curve in one step, until
the power mismatch exceeds \code{param.mismatch_error}. The real roots of the Pad\'e
denominator reveal the location of the nearest singularity on the curve.

\subsection{Hybrid switching and warm start}
When a holomorphic step approaches a singularity --- signalled either by the corrector
failing to converge, or by the holomorphic increment $|\Delta\alpha|$ collapsing --- the
solver switches to predictor--corrector continuation to pass through the fold. To avoid a
slow ``cold'' warm-up, a \textbf{warm starter} sets the first continuation step to a
fraction of the estimated distance to the singularity (and no larger than a fraction of
the last holomorphic step), then builds a quadratic predictor. After the fold is cleared,
control returns to the holomorphic routine. One holomorphic step costs about
$4\times$ a continuation step but typically covers $\approx 8$--$9$ continuation steps of
arc, which is where the speed-up comes from.

\subsection{Continuation details}
The continuation phase uses a pseudo-arclength predictor (quadratic, from the last three
points), a Newton corrector (\emph{Phase I}), and a variable-rescaling corrector
(\emph{Phase II}) for badly-conditioned near-singular arcs whose secant slope exceeds
\code{param.slope_max}. The step length adapts to the corrector iteration count toward a
target of \code{param.targetsteps}. Each $\alpha$ sign change triggers a solution
extraction and a duplicate check; a trace terminates when it returns to its starting
point (within \code{closure} tolerance) or exhausts the alternation cap
\code{param.outnum}.

\subsection{Validated, robustness-hardened configuration}\label{sec:tuned}
This v3 build ships a \emph{validated} default configuration that hardens the tracer
against the numerical failures (wrong-branch drift and non-closing loops near sharp folds)
that can occur with the legacy settings. Four changes matter most, all in
\code{solver_params.m}:
\begin{itemize}[leftmargin=1.4em]
  \item \textbf{Consensus pole estimate} (\code{pole_buses = 2:6}, \code{pole_consensus})
        --- rejects spurious per-bus poles that would place the hand-off badly
        (Section~\ref{sec:pole});
  \item \textbf{Turning-point slope limit} \code{slope_max = 1e4} --- keeps the
        continuation active through steep near-singular arcs instead of handing back too
        early (values of $5\!\times\!10^4$ and above re-break the hard case118 curves);
  \item \textbf{Hand-back hysteresis} \code{handback = 0.5} --- travels far enough past a
        turning point before returning to the holomorphic step that the curve cannot be
        misled onto an adjacent wrong branch;
  \item \textbf{Fold-proximity step caps} \code{init_cap_k = 1500} (caps the first
        continuation step after a hand-back) and \code{max_adapt = true} (an adaptive
        per-step cap driven by the secant-slope change) --- both prevent a large step from
        jumping a sharp apex onto the wrong branch.
\end{itemize}
These defaults are exact on every verifiable case (case3 \dots\ case57mod). Unless you are
experimenting, leave them as shipped; the tables in Section~\ref{sec:params} describe each.

% ================= 7. PARAMETERS =================
\section{Parameter reference}\label{sec:params}

\textbf{All tunables live in \code{solver_params.m}} --- a single function returning the
\code{param} struct that every framework and every parallel worker reads. Edit a line
there and the change applies everywhere. The file is fully commented; the tables below
cover the parameters you are most likely to adjust, with this build's shipped defaults.

\subsection{Continuation / step control}
\begin{center}\small
\begin{longtable}{@{}p{0.24\textwidth}p{0.14\textwidth}p{0.55\textwidth}@{}}
\toprule
\textbf{Parameter} & \textbf{Default} & \textbf{Meaning}\\
\midrule\endhead
\code{sstep}        & \code{1e-5} & Base arc-length step scale for the continuation.\\
\code{targetsteps}  & \code{3}    & Target number of corrector iterations; drives step-size adaptation.\\
\code{imax}         & \code{40}   & Maximum Newton iterations per corrector step.\\
\code{smax}         & \code{20000}& Maximum continuation steps per \code{branch_trace} call.\\
\code{maximumstep}  & \code{8000} & Maximum step-size multiplier (Phase~I ceiling base).\\
\code{maxphase1}    & \code{maximumstep/2} & Maximum step in Phase~I.\\
\code{maxphase2}    & \code{maxphase1/10}  & Maximum step in Phase~II.\\
\code{minimumstep}  & \code{1e-2} & Minimum step-size multiplier.\\
\code{ef}           & \code{1e-8} & Corrector residual tolerance $\lVert F\rVert$.\\
\code{ex}           & \code{1e-7} & Corrector step tolerance $\lVert \Delta x\rVert$.\\
\code{PS}           & \code{300}  & Phase~II (rescaling) duration, in steps.\\
\code{maxsolu}      & \code{100}  & Maximum solutions collected per traced loop.\\
\code{maxskip}      & \code{50}   & Consecutive numerical stalls tolerated before the loop is skipped.\\
\code{dev}          & \code{20}   & Console print interval (display only).\\
\bottomrule
\end{longtable}
\end{center}

\subsection{Holomorphic / hand-off control}
\begin{center}\small
\begin{longtable}{@{}p{0.24\textwidth}p{0.14\textwidth}p{0.55\textwidth}@{}}
\toprule
\textbf{Parameter} & \textbf{Default} & \textbf{Meaning}\\
\midrule\endhead
\code{mismatch_error}& \code{1e-4} & Maximum holomorphic power-mismatch accepted before shrinking the arc.\\
\code{holonum}      & \code{50}   & Maximum Pad\'e steps per outer iteration before forcing continuation.\\
\code{holo_arc_max} & \code{0.01} & Maximum holomorphic arc length per step.\\
\code{outnum}       & \code{500}  & Cap on holomorphic$\leftrightarrow$continuation alternations per curve (a stall burns this cap).\\
\code{slope_max}    & \code{1e4}  & Maximum secant slope at which a turning point is accepted as passed. \code{1e4} is the validated value; $\ge 5\!\times\!10^4$ re-breaks the hard case118 curves.\\
\code{aal_min}      & \code{1e-6} & Minimum homotopy-parameter increment; the unit for \code{handback}.\\
\code{im_tol}       & \code{1e-3} & Imaginary-part tolerance for accepting a Pad\'e root as a real pole.\\
\code{handback}     & \code{0.5}  & Hand-back hysteresis ($\times$\code{aal_min}): how far past a turning point to travel before returning to the holomorphic routine. \code{0.5} is validated across all cases (the legacy build used \code{0.1}).\\
\code{interval}     & \code{1}    & Stride for the turning-point finite-difference test.\\
\bottomrule
\end{longtable}
\end{center}

\subsection{Consensus pole estimation}
\begin{center}\small
\begin{longtable}{@{}p{0.24\textwidth}p{0.14\textwidth}p{0.55\textwidth}@{}}
\toprule
\textbf{Parameter} & \textbf{Default} & \textbf{Meaning}\\
\midrule\endhead
\code{pole_buses}     & \code{2:6}      & Buses sampled for the consensus pole estimate (auto-filtered to $\le$ bus count).\\
\code{pole_consensus}  & \code{true}     & If true, take the smallest pole agreed on by $\ge$ half the sampled buses; else the global minimum.\\
\code{pole_cl_tol}     & \code{0.10}     & Relative magnitude tolerance for clustering ``the same pole'' across buses.\\
\bottomrule
\end{longtable}
\end{center}

\subsection{Fold-proximity and stall control}
These are the robustness controls added in v3 (Section~\ref{sec:tuned}).
\begin{center}\small
\begin{longtable}{@{}p{0.24\textwidth}p{0.14\textwidth}p{0.55\textwidth}@{}}
\toprule
\textbf{Parameter} & \textbf{Default} & \textbf{Meaning}\\
\midrule\endhead
\code{init_cap_k}    & \code{1500} & Caps the first continuation step after a hand-back at \code{init_cap_k}$\times$\code{holo_arc}$/$\code{sstep}, so a fold entry cannot be overshot (\code{Inf} disables).\\
\code{max_adapt}     & \code{true} & Adaptive per-step ceiling: \code{astep} $\le$ \code{maximumstep}$\,\times\min(s_{k-1}/s_k,\,1)$ from the 2-point secant slope $s$ --- shrinks the step approaching a fold, restores it after. \code{false} disables.\\
\code{aal_limit_init}& \code{30}   & Sentinel magnitude meaning ``no Pad\'e pole found yet''.\\
\code{stall_iter_mult}& \code{3}   & Flag a numerical stall when corrector iterations exceed this $\times$\code{targetsteps}.\\
\code{step_adapt_denom}& \code{40} & Step-growth denominator: \code{astep}$\,\times(1+(\text{targetsteps}-\text{iters})/\text{this})$.\\
\bottomrule
\end{longtable}
\end{center}

\subsection{Holomorphic degree}
Set in the drivers / \code{main.m}: \code{degree.act = 15} (active power-series degree),
split as \code{degree.numerator = 8}, \code{degree.denominator = 7} for the Pad\'e
approximant. Changing them changes the Pad\'e order; because this build is pure MATLAB,
the new values take effect immediately (no rebuild needed).

% ================= 8. OUTPUTS =================
\section{Outputs}

\textbf{Console:} per-case progress, solution count, and CPU (serial) or wall (parallel)
time, followed by a summary table.

\textbf{CSV summaries} (written into the run folder):
\begin{center}\small
\begin{tabular}{@{}ll@{}}
\toprule
Driver & CSV file\\
\midrule
\code{run_batch.m}          & \code{results_summary.csv}\\
\code{run_batch_par.m}      & \code{results_par_summary.csv}\\
\code{run_batch_parfeval.m} & \code{results_parfeval_summary.csv}\\
\bottomrule
\end{tabular}
\end{center}
Columns: \code{case_name}, \code{cpu_time_sec} (or wall time), \code{n_solutions},
\code{status} (OK / ERROR), \code{error_msg}.
These CSV files are run outputs and are ignored by the repository \code{.gitignore}.

The full solution set is available in the workspace as \code{Zsave} (each column a
solution in the internal variable ordering) after a run; the solver sign-normalises and
deterministically orders it (\code{canonicalize_solutions.m}) so that serial, parfor and
parfeval produce identical output for the same case. Antipodal pairs $\{x,-x\}$ are
identified as one solution.

% ================= 9. FILE MAP =================
\section{File map}

\textbf{Configuration:} \code{solver_params.m} --- single source of truth for every
tunable constant (Section~\ref{sec:params}).

\textbf{Setup / preprocessing:}
\code{makeYbus} from MATPOWER (Ybus), \code{get_quadr_mtrx.m} (power-balance quadratics),
\code{quadr_matrix.m} (per-bus PV/PQ/slack matrices), \code{Solu.m} (initial point from
\code{runpf}), \code{MakeEllipse.m} (base-ellipse generation), \code{Tune_Fac.m}
(continuation scaling), \code{main.m} (reference preprocessing script; its logic is
inlined into the drivers).

\textbf{Core solver (shared hot path):}
\code{hybrid_traceloops_4_no_mex.m} (outer trace-loop orchestrator; loads \code{param}
from \code{solver_params.m}, primes the cached operators), \code{holomorphic_hybrid_5_no_mex.m} (one
holomorphic step + consensus pole estimate), \code{branch_trace_hybrid_4_no_mex.m}
(predictor--corrector continuation), \code{Resolve_no_mex.m} (Newton corrector),
\code{holomorphic_cont_tri.m} (Taylor-coefficient triangular solves),
\code{holomorphic_para_sp.m} (sparse $H$ builder), \code{Pade_Apprxmt.m} (Pad\'e
construction), \code{update_V_Pade.m} (Pad\'e evaluation along the arc),
\code{MakeJacobianD.m} (cached sparse Jacobian), \code{quad_form_vals.m} (cached
constraint-residual evaluator), \code{MakeJacobian.m} (scalar Jacobian, used by
\code{Tune_Fac}).

\textbf{Post-processing:} \code{canonicalize_solutions.m} (sign-normalise + deterministic
ordering).

\textbf{Parallel housekeeping:} \code{cleanup_parallel_artifacts.m} (clears per-worker
caches, optionally closes the pool to avoid a Windows/R2022a \code{-batch} shutdown
crash).

\textbf{Drivers:}
serial --- \code{runVBook_hybrid_2023.m}, \code{run_batch.m};
parfor --- \code{runVBook_hybrid_parallel.m}, \code{trace_equation_worker.m},
\code{run_batch_par.m};
parfeval --- \code{runVBook_hybrid_parfeval.m}, \code{trace_equation_worker_const.m},
\code{trace_equation_worker.m}, \code{run_batch_parfeval.m};
convenience --- \code{run_merged_case.m} (one case, any mode).

\textbf{Case data:} \code{caseXXX.m} MATPOWER-format test systems.

% ================= 10. PORTABILITY =================
\section{Portability and performance}

This is the pure-MATLAB build: no compilation, no MATLAB Coder, and no C compiler --- it
runs unchanged on Windows, Linux and macOS. Every hot-path routine is a plain \code{.m}
file (e.g.\ \code{Pade_Apprxmt.m}, \code{holomorphic_cont_tri.m}), so editing and
debugging require nothing beyond MATLAB itself.

If you need more speed on Windows x64, the companion release \code{HEBCPF_MEX_v3.202606}
ships two compiled MEX functions on the hot path (\code{Pade_Apprxmt} and
\code{holomorphic_cont_tri}) and runs the serial trace roughly $2\times$ faster, with
numerically identical results (verified to $\sim 10^{-13}$). The two builds are otherwise
interchangeable: you can develop here and deploy there, or vice versa.

% ================= 11. TROUBLESHOOTING =================
\section{Troubleshooting}

\begin{description}[leftmargin=0em,style=nextline]
\item[\code{Undefined function 'runpf'}] MATPOWER is not on the MATLAB path. Install and
      \code{addpath} it.
\item[\code{Undefined function 'makeYbus'}] MATPOWER is not on the MATLAB path, or the
      MATPOWER \code{lib} folder is missing from the path.
\item[\code{Undefined function 'caseXXX'}] Run \code{addpath(pwd)} from this folder; the
      case data files live here.
\item[OSQP warning during \code{runpf}] Harmless MATPOWER solver-probe message;
      \code{runpf} still converges (equation error $\sim 0$).
\item[Parallel driver appears to hang at first run] The local \code{parpool} is starting;
      this one-time cost is excluded from the reported timing.
\item[Solution count differs from the expected value] Check the load factor is 1 and that
      the MATPOWER initial solve (\code{runpf}) reported success. Note that the
      pole-estimation, \code{handback} and \code{slope_max} settings affect \emph{tracing
      efficiency and numerical robustness}, not the final solution set on well-behaved
      cases.
\end{description}

% ================= 12. LICENSE =================
\section{License and third-party notices}

HEBCPF is distributed under the BSD 3-Clause License; see the top-level
\code{LICENSE} file. Some bundled \code{caseXXX.m} files are copied from, derived from,
or modified from MATPOWER and other public test-system sources. Preserve the per-file
attribution notices when redistributing modified versions, and see the top-level
\code{NOTICE} file for third-party attribution details.

% ================= 12. CITATION =================
\section{Citing this solver}\label{sec:cite}

\textbf{If you use results produced by this solver, please cite the HEBC paper:}
\begin{quote}
D.\ Wu and B.\ Wang, ``Holomorphic Embedding Based Continuation Method for Identifying
Multiple Power Flow Solutions,'' \emph{IEEE Access}, vol.~7, pp.~86843--86853, 2019.
doi:\,\doi{10.1109/ACCESS.2019.2925384}.
\end{quote}

BibTeX:
\begin{lstlisting}[style=sh]
@article{wu2019hebc,
  author  = {Wu, Dan and Wang, Bin},
  title   = {Holomorphic Embedding Based Continuation Method for
             Identifying Multiple Power Flow Solutions},
  journal = {IEEE Access},
  volume  = {7},
  pages   = {86843--86853},
  year    = {2019},
  doi     = {10.1109/ACCESS.2019.2925384}
}
\end{lstlisting}

For citing the software package itself, see the top-level \code{CITATION.cff} file.

\subsection*{Further reading}
The method rests on, and is extended by, the following works:

\begin{enumerate}[leftmargin=1.6em]
\item B.\ C.\ Lesieutre and D.\ Wu, ``An Efficient Method to Locate All the Load Flow
      Solutions --- Revisited,'' in \emph{2015 53rd Annual Allerton Conference on
      Communication, Control, and Computing (Allerton)}, Monticello, IL, USA, pp.~381--388,
      IEEE, 2015. doi:\,\doi{10.1109/ALLERTON.2015.7447029}. \emph{(The elliptical
      branch-tracing formulation that underlies this solver.)}
\item D.\ Wu, ``Algebraic Set Preserving Mappings for Electric Power Grid Models and Its
      Applications,'' Ph.D.\ dissertation, University of Wisconsin--Madison, 2017.
      \emph{(The dissertation that formalises the elliptical / algebraic-set-preserving
      formulation and specifies the predictor--corrector continuation machinery --- the
      Phase~I/II correctors, step-length control, and solution/termination logic --- that
      this solver implements.)}
\item D.\ Wu, D.\ K.\ Molzahn, B.\ C.\ Lesieutre, and K.\ Dvijotham, ``A Deterministic
      Method to Identify Multiple Local Extrema for the AC Optimal Power Flow Problem,''
      \emph{IEEE Transactions on Power Systems}, vol.~33, no.~1, pp.~654--668, Jan.\ 2018.
      doi:\,\doi{10.1109/TPWRS.2017.2707925}.
      \emph{(Extends the tracing method to AC optimal power flow.)}
\item B.\ Wang, D.\ Wu, X.\ Su, K.\ Sun, and L.\ Xie, ``A Long-Term Voltage Stability
      Margin Index Based on Multiple Real Power Flow Solutions,'' in \emph{2023 IEEE Power
      \& Energy Society General Meeting (PESGM)}, Orlando, FL, USA, pp.~1--5, IEEE,
      Jul.\ 2023. doi:\,\doi{10.1109/PESGM52003.2023.10252776}.
      \emph{(Defines a voltage-stability margin as the
      minimum voltage-space distance from the high-voltage solution to all other real
      power-flow solutions --- a direct application of this solver's output.)}
\item D.\ Wu and B.\ Wang, ``Influence of Load Models on Equilibria, Stability and
      Algebraic Manifolds of Power System Differential-Algebraic System,'' in \emph{2019
      57th Annual Allerton Conference on Communication, Control, and Computing
      (Allerton)}, Monticello, IL, USA, IEEE, 2019.
      doi:\,\doi{10.1109/ALLERTON.2019.8919649}. \emph{(How load models reshape the
      solution manifold and its equilibria.)}
\item D.\ Wu and B.\ Wang, ``Collection of Numerous Power Flow Solutions of Standard IEEE
      Test Systems,'' \emph{IEEE Dataport}, Jul.\ 2019.
      doi:\,\doi{10.21227/24bh-hj72}. \emph{(The published dataset of reference solution
      sets for the standard IEEE test systems --- useful for validating this solver's
      output.)}
\end{enumerate}

\vfill
\begin{center}\small\itshape
HEBCPF --- MATLAB v3.202606. This document accompanies the solver release and may be
redistributed with it.
\end{center}

\end{document}
